% ═══════════════════════════════════════════════════════════════════════
%  Stratégie 2 --- XAUUSD niveaux x10. Cette section pose le CADRE de la
%  réconciliation ; les mesures QC et MT5 arrivent au lot suivant et ne
%  sont représentées que par des macros « n.d. ».
% ═══════════════════════════════════════════════════════════════════════
\section{Recherche contre exécution : le cadre de réconciliation}
\label{sec:x10_reconciliation}

\lettrine[lines=2, findent=2pt, nindent=0pt]{U}{ne} stratégie mesurée par un
seul moteur n'a été mesurée qu'une fois. Ce chapitre décrit le dispositif qui
confronte trois implémentations indépendantes de la même spécification, ce
qu'il peut établir, et ce qu'il ne peut pas.

\ifdefstring{\XTenVerdictProvisoire}{1}{%
\begin{warningbox}[title=\textbf{\color{white}\textsf{Mesures QuantConnect et MetaTrader~5 : en cours d'intégration}}]
Les deux portages d'exécution existent et leurs campagnes sont en cours. Leurs
résultats ne sont pas publiés ici, et aucun chiffre d'exécution n'est avancé~:
l'espérance QuantConnect vaut \code{\XTenQCExpectancyR}, l'espérance
MetaTrader~5 \code{\XTenMTFiveExpectancyR}, et les taux d'appariement
\code{\XTenQCMatchRate} et \code{\XTenMTFiveMatchRate}. Ces valeurs seront
remplies par le même générateur que le reste du rapport, à partir des
fichiers de résultats des deux moteurs, sans intervention manuelle.

\textbf{Cette attente ne suspend pas le verdict.} La table de décision compte
un critère non mesuré comme un échec~: le critère MetaTrader~5 est donc déjà
compté en échec, et une mesure ultérieure ne peut que confirmer ou aggraver le
constat, jamais l'annuler --- six critères mesurés échouent par ailleurs.
\end{warningbox}
}{}

\subsection{Ce que chaque moteur simule}

\begin{center}
\renewcommand{\arraystretch}{1.35}
\begin{tabular}{p{3.0cm}p{11.0cm}}
\toprule
\textbf{Moteur vectoriel} & Évalue la stratégie sur des tableaux entiers,
d'un bloc. Il permet d'explorer une grille de vingt-sept configurations en
quelques secondes, ce qui est indispensable pour mesurer un plateau plutôt
qu'un point. Il ne simule pas le passage du temps~: il le calcule. C'est le
moteur de \emph{recherche}. \\
\textbf{QuantConnect} & Rejoue la stratégie barre par barre, sur les
\emph{mêmes} données que le moteur vectoriel, avec un livre d'ordres, un
gestionnaire de portefeuille et des horaires de marché. Il sert à démontrer
que la stratégie est \emph{implémentable en temps réel}, et que le moteur
vectoriel n'a pas triché sur la causalité. \\
\textbf{MetaTrader~5} & Rejoue la stratégie sur le flux du \emph{courtier},
qui n'est pas celui des deux autres, avec son spread, ses arrondis de lots et
ses règles de marge. C'est le moteur qui passerait les ordres~: ses chiffres
sont ceux qui feraient foi pour le compte. \\
\bottomrule
\end{tabular}
\end{center}

Les deux premiers partagent les mêmes barres~; le troisième non, et cette
asymétrie n'est pas un défaut à corriger. L'historique de l'or chez le
courtier commence fin 2022, ce qui limite structurellement la comparaison à
un peu plus de trois des sept années d'échantillon --- et ce sont les années
où le coût pèse le moins lourd, donc les plus favorables. Les régimes de
2019 à 2022 ne sont \emph{pas} vérifiables sur le moteur d'exécution.

\subsection{Protocole de réconciliation événement par événement}

Les trois moteurs émettent le même fichier de trace~: une ligne par événement,
vingt-et-une colonnes dans un ordre fixé, horodatage en temps universel,
arrondis imposés. L'automate n'a le droit d'émettre que six types
d'événements --- armement, cassure, débordement, entrée, sortie, annulation ---
et de changer d'état sans émettre la ligne correspondante n'est pas permis.

\begin{insightbox}
\textbf{Comparer des P\&L ne sert à rien~; comparer des événements localise la
cause.} Deux moteurs qui affichent le même résultat final peuvent être en
désaccord sur la moitié de leurs transactions, les erreurs se compensant. Deux
moteurs qui affichent des résultats différents peuvent n'avoir qu'un seul
désaccord, dans le calcul d'un indicateur. Seule la trace événementielle
distingue les deux cas.
\end{insightbox}

\subsection{L'échelle de lecture à cinq barreaux}

La comparaison descend barreau par barreau et s'arrête au premier qui casse~:
un écart au barreau \emph{n} rend tous les barreaux au-delà ininterprétables.

\begin{center}
\renewcommand{\arraystretch}{1.35}
\begin{tabular}{c p{4.6cm} p{7.6cm}}
\toprule
\textbf{Barreau} & \textbf{Objet comparé} & \textbf{Si ça casse, la cause est} \\
\midrule
1 & Les barres de cinq minutes & Bornes de séance, fuseau horaire,
calendrier \\
2 & Les indicateurs~: \code{ATR}, vitesse, momentum, accélération,
\code{VWAP}, moyenne horaire, dollar & Lissage, amorçage, causalité des
séries horaires \\
3 & Les événements d'automate & Seuils, largeur de zone, maintien,
temporisation après sortie \\
4 & Les entrées & Garde-fou gain/risque, fenêtre horaire, unicité de
position \\
5 & Les sorties et le P\&L & Côté acheteur/vendeur, double contact, écart de
cotation, lots, coûts \\
\bottomrule
\end{tabular}
\end{center}

\subsection{Écarts structurels attendus, et tolérances}

Les tolérances ont été fixées dans la spécification, \emph{avant} la première
comparaison, et elles ne sont pas les mêmes selon les moteurs confrontés.

% <<< LITERAL-OK : tolérances de réconciliation et seuil de divergence non
% attribuée. Ce sont des SEUILS gelés dans la spécification avant toute
% comparaison, pas des mesures : aucun fichier de résultats ne les porte.
\begin{center}
\renewcommand{\arraystretch}{1.35}
\begin{tabular}{p{5.2cm} p{2.6cm} p{5.6cm}}
\toprule
\textbf{Comparaison} & \textbf{Données} & \textbf{Appariement attendu} \\
\midrule
Python contre QuantConnect & identiques & au moins 98~\% des entrées \\
Python contre MetaTrader~5, sur données du courtier & identiques & au moins
95~\% des entrées \\
Python ou QuantConnect contre MetaTrader~5 & différentes & au moins 70~\%,
l'écart devant être \emph{attribué} \\
\bottomrule
\end{tabular}
\end{center}

\begin{warningbox}
\textbf{Un écart n'est un échec que s'il reste inexpliqué.} Viser l'égalité
avec le moteur du courtier serait le signe qu'on a idéalisé le backtest, pas
qu'on a réconcilié quoi que ce soit~: les flux diffèrent, les niveaux frôlés
d'un côté ne le sont pas de l'autre, et le prix médian n'est pas le prix
acheteur. La règle opposable est différente~: une divergence \textbf{non
attribuée} portant sur plus de 5~\% des transactions \emph{bloque} la lecture
hors échantillon.
\end{warningbox}
% >>> LITERAL-OK

Trois sources d'écart sont attendues et documentées d'avance.

\begin{description}
  \item[L'amorçage des indicateurs.] L'\code{ATR} natif de MetaTrader~5 ne
  s'amorce pas comme celui de la spécification, et les deux moteurs
  d'exécution doivent précharger au moins 51~barres horaires avant la première
  barre évaluée. Sans cela, le barreau~2 casse d'entrée.
  \item[Le modèle de remplissage du testeur.] Le testeur tourne en mode
  « OHLC M1 », qui \emph{interpole} les remplissages à l'intérieur de la
  minute. Les chiffres MetaTrader~5 sont donc structurellement un
  \textbf{majorant}, et doivent être lus comme tels.
  \item[Le panier dollar multi-symboles.] La récupération simultanée de
  quatre symboles n'est pas garantie dans le testeur. La spécification prévoit
  le cas~: une jambe absente depuis plus de cinq barres rend le contexte
  dollar \emph{neutre}, et l'occurrence est comptée.
\end{description}

\subsection{Ce que la réconciliation ne pourra pas établir}

Même achevée, elle ne dira pas si la stratégie gagne de l'argent~: cette
question est tranchée par la recherche, et elle l'est négativement. Elle dira
si les trois moteurs décrivent le même objet. C'est une condition nécessaire
à un déploiement, jamais une condition suffisante --- et, dans ce dossier,
c'est une condition dont la vérification ne changera pas la décision.
